In Table~\ref{tab:crypto}, some global cryptographic commands are shown.
\orbitertableofcommands{cryptography_1}
Some of these commands require a finite field 
and appear as a finite field activity, 
see Table~\ref{tab:crypto2}.
\orbitertableofcommands{cryptography_2}
For instance,

\sourcecode{EC_add}

adds the point $(1,4)$ on the curve $y^2 = x^3 + x + 3 \mod 11$ to itself.
The command 

\sourcecode{EC_cyclic_subgroup}

computes the cyclic subgroup generated by the point $(1,4)$ 
on the curve $y^2 = x^3 + x + 3 \mod 11.$
The command

\sourcecode{EC_points_199}

computes all points on the curve $y^2 = x^3 + 5x + 7 \mod 199$ 
and produces a bitmap drawing of the points in the affine plane:
\orbiteroutput{GRAPHICS/WRAPPERS/EC_5_7_q199_points_xy_draw}
Both the $x$-axis and the $y$-axis 
are indexed by the field elements from 0 to 198.


\bigskip

The command

\sourcecode{EC_Koblitz_encoding}

encode the message ``DEADBEEF'' on the curve  $y^2 = x^3 + 5x + 7 \mod 199$ 
using the base point $(147,164)$ and the secret key $67.$
The $i$th input character is encoded as two points 
$(R_i,T_i)$ on the curve using the Elgamal scheme.
A random round key is generated for each plaintext symbol. 
As seen in this example, the \verb'-seed' 
command can be used to seed the random number generator with 
an arbitrary integer (here 17).

\bigskip


The command

\sourcecode{EC_bsgs}

performs a baby-step-giant-step brute force attack on the ciphertext sequence
\begin{align*}
R_i &= (172,158),(45,195),(50,22),(10,103),(55,33),\\
&(50,22),(145,105),(31,74),(73,155),(67,60),(25,6),
\end{align*} 
using the base point $(147,164)$ on the curve  $y^2 = x^3 + 5x + 7 \mod 199,$ 
assuming a group order of $212.$ 
The command 

\sourcecode{EC_bsgs_decode}

decodes the ciphertext sequence
\begin{align*}
T_i &= (127,188),(51,141),(85,29),(106,90),(41,105),(179,71),\\
&(171,2),(16,197),(183,72),(27,129),(37,10),
\end{align*} 
assuming round keys 
$$
k_i = 50,179,169,13,153,169,115,116,188,110,176,
$$
using the base point $(147,164)$ on the curve  $y^2 = x^3 + 5x + 7 \mod 199,$ 
and assuming a group order of $212.$ 

\bigskip


The next sequence of examples discusses the NTRU cryptosystem 
(cf. Example 7.53 in~\cite{Hoffstein2014}).
In the example, we choose the parameters of the cryptosystem to be $(N,p,q,d) = (7,41,3,2).$
Orbiter uses the following convention for polynomials over a finite field $\bbF_q$: 
The coefficients of $A(X) = a_0 + a_1X + \cdots + a_dX^d$ 
are listed as a sequence, starting with the constant term and ending with the leading coefficient.
The cryptosystem requires coefficients $a_i$ in the range $-\frac{p}{2} \le a_i \le \frac{p}{2}.$ 
So, in an extension to the conventionas for field elements in $\bbF_q$, 
Orbiter allows negative coefficients as well. 
The assumption is that $q$ is prime and 
negative coefficients are considered modulo $q.$
In the example, Alice picks the private polynomials 
$f(x) = x^6-x^4+x^3+x^2-1$ (with $d+1$ coefficients equal to 
plus one and $d$ coefficients equal to minus one) and 
$g(x) = x^6+x^4-x^2-x$ with $d$ coefficients 
plus one and $d$ coefficients minus one.
We also need the polynomial $x^N-1$. 
The makefile commands 

\sourcecodevariable{NTRU_N}


\sourcecodevariable{NTRU_P}


\sourcecodevariable{NTRU_Q}

\sourcecodevariable{NTRU_D}


\sourcecodevariable{NTRU_XN1}


\sourcecodevariable{ALICE_PRIVATE_F}


\sourcecodevariable{ALICE_PRIVATE_G}

are used to set up the appropriate variables according to these  choices. 

\bigskip


Regarding the NTRU set-up, Alice needs to compute her private keys 
$F_p(x)$ and $F_q(x).$ These two polynomials are defined as follows:
\begin{enumerate}
\item
	$F_p(x)$ is the inverse of $f(x)$ in $\bbF_p[x]/(x^n-1),$ 
\item
$F_q(x)$ the inverse of $f(x)$ in $\bbF_q[x]/(x^n-1).$
\end{enumerate}
To this end, we can use the \verb'extended_gcd_for_polynomials' 
command from Table~\ref{tab:nt}.
The following two makefile commands do the job: 

\sourcecode{NTRU_Alice1}


\sourcecodevariable{ALICE_PRIVATE_FQ}


\sourcecode{NTRU_Alice2}

The resulting polynomials 
(indicated as comments by means of the  \verb'#' symbol) 
are again encoded as makefile variables. 

\sourcecodevariable{ALICE_PRIVATE_FP}

There is a chance that the polynommial 
$f(x)$ does not have an inverse in either $\bbF_p[x]$ 
or in $\bbF_q[x].$ 
In that case, Alice simply chooses a different 
polynomial $f(x)$ and tries again.
Alice can now compute her public key:

\sourcecode{NTRU_Alice_public_key}


\sourcecodevariable{ALICE_PUBLIC_KEY}

The public key is assigned to the makefile variable \verb'ALICE_PUBLIC_KEY'.
Now, Bob chooses his message to Alice and his one-time-key. 
The message must be the center lift of a polynomial in $\bbF_p[x].$
The round-key must have exactly $d$ coefficients one 
and $d$ coefficients $-1$ (rest zeroes).

\sourcecodevariable{BOB_MESSAGE}


\sourcecodevariable{BOB_ONE_TIME_KEY}

The encryption proceeds using the \verb'NTRU_encrypt' command, 
and the result is stored in the makefile variable \verb'BOB_ENCRYPT':

\sourcecode{NTRU_encrypt}


\sourcecodevariable{BOB_ENCRYPT}

Decryption is done in five steps. 

\sourcecode{NTRU_decrypt1}

\sourcecodevariable{ALICE_C1}


\sourcecode{NTRU_decrypt2}


\sourcecodevariable{ALICE_C2}


\sourcecode{NTRU_decrypt3}


\sourcecodevariable{ALICE_C3}


\sourcecode{NTRU_decrypt4}


\sourcecodevariable{ALICE_C4}


\sourcecode{NTRU_decrypt5}

Decryption produces Bob's message to Alice.

\bigskip

%ToDo:

%\begin{itemize}
%\item
%RSA
%\item
%sqrt mod
%\item
%quadratic sieve
%\item
%pseudoprimes
%\end{itemize}

%\bigskip

